\documentclass[11pt,a4paper]{article}
\usepackage[utf8]{inputenc}
\usepackage[T1]{fontenc}
\usepackage[french]{babel}
\usepackage{amsmath}
\usepackage{booktabs}
\usepackage{enumitem}
\usepackage[a4paper,margin=1.7cm]{geometry}
\pagestyle{empty}
\setlength{\parindent}{0pt}
\newcommand{\champ}[1]{\makebox[#1]{\dotfill}}
\begin{document}

\begin{center}{\LARGE\bfseries Le logarithme décimal}\\[4pt]{\large Deux équipes au travail : papier contre ordinateur}\\[8pt]{\small Feuille de route enseignant --- à compiler avec pdflatex (TeXStudio : touche F5)}\end{center}

{\bfseries Classe :} \champ{4.6cm}\hfill {\bfseries Date :} \champ{3.6cm}\par\smallskip
{\bfseries Effectif :} deux équipes de 8 élèves\quad\quad{\bfseries Exploration :} 50 minutes par rotation, puis inversion des groupes
\section*{Principe}
Deux équipes de 8 élèves explorent pendant 50 minutes les six domaines de la feuille de découverte, chacune avec ses propres moyens. L'équipe \emph{papier} n'a qu'une feuille, des stylos et des cerveaux : elle raisonne, calcule à la main et débat. L'équipe \emph{ordinateur} mène une recherche guidée sur le web et note ses trouvailles avec leurs sources. Puis on intervertit les groupes : chaque équipe retrouve les mêmes six domaines sous l'autre angle, et compare ce que l'autre support permet de découvrir.
\section*{Déroulé}
\begin{tabular}{@{}llp{0.52\linewidth}@{}}
\toprule
Horaires & Durée & Séquence \\
\midrule
0 h 00 --- 0 h 10 & 10 min & Installation, constitution des équipes A et B, lecture des règles \\
0 h 10 --- 1 h 00 & 50 min & \textbf{Rotation 1} --- équipe A : papier \quad/\quad équipe B : ordinateur \\
1 h 00 --- 1 h 05 & 5 min & Inversion des groupes : chacun change de support \\
1 h 05 --- 1 h 55 & 50 min & \textbf{Rotation 2} --- équipe A : ordinateur \quad/\quad équipe B : papier \\
1 h 55 --- 2 h 00 & 5 min & Bilan collectif : mise en commun avec la grille \\
\bottomrule
\end{tabular}
\section*{Composition des équipes}
\begin{tabular}{@{}p{0.46\linewidth}@{\hspace{0.05\linewidth}}p{0.46\linewidth}@{}}
\textbf{Équipe A} (papier, puis ordinateur) & \textbf{Équipe B} (ordinateur, puis papier) \\
\midrule
1. \champ{0.8\linewidth} & 1. \champ{0.8\linewidth} \\
2. \champ{0.8\linewidth} & 2. \champ{0.8\linewidth} \\
3. \champ{0.8\linewidth} & 3. \champ{0.8\linewidth} \\
4. \champ{0.8\linewidth} & 4. \champ{0.8\linewidth} \\
5. \champ{0.8\linewidth} & 5. \champ{0.8\linewidth} \\
6. \champ{0.8\linewidth} & 6. \champ{0.8\linewidth} \\
7. \champ{0.8\linewidth} & 7. \champ{0.8\linewidth} \\
8. \champ{0.8\linewidth} & 8. \champ{0.8\linewidth} \\
\end{tabular}
\section*{Matériel et impressions}
\begin{itemize}[leftmargin=1.4em,itemsep=2pt]
\item \textbf{Fiche équipe papier :} 16 exemplaires (8 par rotation --- les élèves écrivent dessus).
\item \textbf{Fiche équipe ordinateur :} 16 exemplaires, et 4 postes avec accès web (un par binôme).
\item Équipe papier : stylos, feuilles de brouillon à volonté. \textbf{Interdits :} téléphone, calculatrice, ordinateur.
\item Grille de mise en commun (dernière page) : une par équipe, ou une copie agrandie au tableau.
\end{itemize}
\section*{Pendant les rotations}
\begin{itemize}[leftmargin=1.4em,itemsep=2pt]
\item Répondez aux questions des élèves par des questions : l'exploration leur appartient.
\item Côté papier, l'interdit du téléphone est total ; encouragez le débat à voix haute et les calculs détaillés.
\item Côté ordinateur, rappelez la règle : une trouvaille sans source n'est pas une trouvaille.
\item À la fin de chaque rotation : un rapporteur pour l'équipe papier, un binôme porte-parole pour l'équipe ordinateur.
\end{itemize}
\section*{Après la séance}
Les élèves reportent leurs découvertes dans le carnet de bord du chapitre interactif (questions d'approfondissement), téléchargent leur carnet chiffré, et vous l'analysez avec le lecteur enseignant.
\newpage
\begin{center}{\Large\bfseries Équipe PAPIER --- explorer sans machines}\\[4pt]{\small Le logarithme décimal : d'où il vient, à quoi il sert --- exploration de 50 minutes}\end{center}
{\small Classe : \champ{3.2cm}\hfill Prénom NOM : \champ{4.6cm}}
\subsection*{Règles du jeu}
\begin{itemize}[leftmargin=1.4em,itemsep=2pt]
\item Vos seuls outils : cette feuille, des stylos, vos cerveaux. Téléphone, calculatrice et ordinateur sont interdits.
\item Explorez les six domaines, dans l'ordre de votre choix. Répondez en une à trois phrases, en montrant vos calculs.
\item En fin de rotation, l'équipe choisit un rapporteur : il présentera deux découvertes au reste de la classe.
\end{itemize}
{\bfseries Boîte à outils :} $\log 2 \approx 0{,}301$, $\log 3 \approx 0{,}477$, $\log 5 \approx 0{,}699$ ; $10^{k}$ possède $k+1$ chiffres ; $+10$ dB $= \times 10$ en puissance sonore ; $+1$ de magnitude $\approx \times 10$ en amplitude.\par\medskip
\subsection*{Domaine 1 --- Histoire : pourquoi le logarithme a-t-il été inventé ?}
\textbf{1.} Sans aucun document : pourquoi des mathématiciens du début du XVII\textsuperscript{e} siècle auraient-ils besoin d'un nouvel outil ? Que font, à la main, les astronomes, les navigateurs et les marchands de l'époque ?
\par\smallskip
\noindent\makebox[\linewidth]{\dotfill}\par\vspace{1em}
\noindent\makebox[\linewidth]{\dotfill}\par\vspace{1em}
\noindent\makebox[\linewidth]{\dotfill}\par\vspace{1em}
\par\medskip
\textbf{2.} Expliquez avec vos mots comment remplacer la multiplication $1\,000 \times 100\,000$ par une simple addition. En général, que devient la multiplication $a \times b$ ?
\par\smallskip
\noindent\makebox[\linewidth]{\dotfill}\par\vspace{1em}
\noindent\makebox[\linewidth]{\dotfill}\par\vspace{1em}
\noindent\makebox[\linewidth]{\dotfill}\par\vspace{1em}
\par\medskip
\subsection*{Domaine 2 --- Le principe : transformer les produits en sommes}
\textbf{1.} Complétez : $2^{4} \times 2^{3} = 2^{\ldots}$ ; $10^{2} \times 10^{5} = 10^{\ldots}$ ; $10^{7} \times 10^{-3} = 10^{\ldots}$.
\par\smallskip
\noindent\makebox[\linewidth]{\dotfill}\par\vspace{1em}
\noindent\makebox[\linewidth]{\dotfill}\par\vspace{1em}
\noindent\makebox[\linewidth]{\dotfill}\par\vspace{1em}
\par\medskip
\textbf{2.} On sait que $2^{10} = 1\,024 \approx 10^{3}$. Sans calculatrice, encadrez $2^{20}$ entre deux puissances de 10. Combien de chiffres possède-t-il ?
\par\smallskip
\noindent\makebox[\linewidth]{\dotfill}\par\vspace{1em}
\noindent\makebox[\linewidth]{\dotfill}\par\vspace{1em}
\noindent\makebox[\linewidth]{\dotfill}\par\vspace{1em}
\par\medskip
\textbf{3.} Une multiplication devient une addition : rédigez la règle générale en une phrase, puis testez-la sur $10^{4} \times 10^{6}$.
\par\smallskip
\noindent\makebox[\linewidth]{\dotfill}\par\vspace{1em}
\noindent\makebox[\linewidth]{\dotfill}\par\vspace{1em}
\noindent\makebox[\linewidth]{\dotfill}\par\vspace{1em}
\par\medskip
\subsection*{Domaine 3 --- Lire un exposant sans machine}
\textbf{1.} Complétez : $\log 1 = \,?$ ; $\log 10 = \,?$ ; $\log 1000 = \,?$ ; $\log 0{,}01 = \,?$.
\par\smallskip
\noindent\makebox[\linewidth]{\dotfill}\par\vspace{1em}
\noindent\makebox[\linewidth]{\dotfill}\par\vspace{1em}
\noindent\makebox[\linewidth]{\dotfill}\par\vspace{1em}
\par\medskip
\textbf{2.} Expliquez pourquoi la partie entière de $\log b$ vaut (nombre de chiffres de $b$) moins un.
\par\smallskip
\noindent\makebox[\linewidth]{\dotfill}\par\vspace{1em}
\noindent\makebox[\linewidth]{\dotfill}\par\vspace{1em}
\noindent\makebox[\linewidth]{\dotfill}\par\vspace{1em}
\par\medskip
\textbf{3.} À l'aide de $\log 2 \approx 0{,}301$, trouvez le nombre de chiffres de $2^{100}$. Décrivez votre raisonnement pas à pas.
\par\smallskip
\noindent\makebox[\linewidth]{\dotfill}\par\vspace{1em}
\noindent\makebox[\linewidth]{\dotfill}\par\vspace{1em}
\noindent\makebox[\linewidth]{\dotfill}\par\vspace{1em}
\par\medskip
\subsection*{Domaine 4 --- Le logarithme dans la vie quotidienne}
\textbf{1.} Le pH d'un citron vaut environ 2,5 ; celui de l'eau pure vaut 7. Que signifie un écart de 1 sur l'échelle de pH, en termes de concentration ? Et un écart de 4,5 ?
\par\smallskip
\noindent\makebox[\linewidth]{\dotfill}\par\vspace{1em}
\noindent\makebox[\linewidth]{\dotfill}\par\vspace{1em}
\noindent\makebox[\linewidth]{\dotfill}\par\vspace{1em}
\par\medskip
\textbf{2.} Une conversation chuchotée atteint 30 dB, un concert 90 dB. Si $+10$ dB multiplie la puissance sonore par 10, par combien est-elle multipliée entre les deux ?
\par\smallskip
\noindent\makebox[\linewidth]{\dotfill}\par\vspace{1em}
\noindent\makebox[\linewidth]{\dotfill}\par\vspace{1em}
\noindent\makebox[\linewidth]{\dotfill}\par\vspace{1em}
\par\medskip
\textbf{3.} Citez encore une situation de la vie quotidienne où se cache une échelle logarithmique, et justifiez en une phrase.
\par\smallskip
\noindent\makebox[\linewidth]{\dotfill}\par\vspace{1em}
\noindent\makebox[\linewidth]{\dotfill}\par\vspace{1em}
\noindent\makebox[\linewidth]{\dotfill}\par\vspace{1em}
\par\medskip
\subsection*{Domaine 5 --- Sciences et industrie : compresser les échelles}
\textbf{1.} De la taille d'un atome ($10^{-10}$ m) à celle de l'univers observable ($10^{26}$ m), combien d'ordres de grandeur séparent les extrêmes ? Pourquoi une échelle logarithmique les fait-elle tenir sur une seule page ?
\par\smallskip
\noindent\makebox[\linewidth]{\dotfill}\par\vspace{1em}
\noindent\makebox[\linewidth]{\dotfill}\par\vspace{1em}
\noindent\makebox[\linewidth]{\dotfill}\par\vspace{1em}
\par\medskip
\textbf{2.} Une épidémie double chaque semaine. Décrivez (ou dessinez) l'allure de la courbe du nombre de cas en échelle ordinaire, puis en échelle log : qu'observe-t-on ?
\par\smallskip
\noindent\makebox[\linewidth]{\dotfill}\par\vspace{1em}
\noindent\makebox[\linewidth]{\dotfill}\par\vspace{1em}
\noindent\makebox[\linewidth]{\dotfill}\par\vspace{1em}
\par\medskip
\textbf{3.} Pourquoi les tolérances des composants électroniques s'expriment-elles souvent en décibels ?
\par\smallskip
\noindent\makebox[\linewidth]{\dotfill}\par\vspace{1em}
\noindent\makebox[\linewidth]{\dotfill}\par\vspace{1em}
\noindent\makebox[\linewidth]{\dotfill}\par\vspace{1em}
\par\medskip
\subsection*{Domaine 6 --- Informatique : chercher en un minimum d'étapes}
\textbf{1.} Un dictionnaire contient $2^{20}$ pages. Montrez que vingt questions « avant ou après ? » suffisent pour trouver n'importe quel mot (raisonnement par dichotomie).
\par\smallskip
\noindent\makebox[\linewidth]{\dotfill}\par\vspace{1em}
\noindent\makebox[\linewidth]{\dotfill}\par\vspace{1em}
\noindent\makebox[\linewidth]{\dotfill}\par\vspace{1em}
\par\medskip
\textbf{2.} Un algorithme « en $\log n$ » double-t-il le nombre d'étapes quand on double les données ? Expliquez en comparant à un algorithme qui lit tout.
\par\smallskip
\noindent\makebox[\linewidth]{\dotfill}\par\vspace{1em}
\noindent\makebox[\linewidth]{\dotfill}\par\vspace{1em}
\noindent\makebox[\linewidth]{\dotfill}\par\vspace{1em}
\par\medskip
\newpage
\begin{center}{\Large\bfseries Équipe ORDINATEUR --- recherche guidée sur le web}\\[4pt]{\small Le logarithme décimal : d'où il vient, à quoi il sert --- exploration de 50 minutes}\end{center}
{\small Classe : \champ{3.2cm}\hfill Prénom NOM : \champ{3.2cm}\hfill Binôme : \champ{2.2cm}}
\subsection*{Règles du jeu}
\begin{itemize}[leftmargin=1.4em,itemsep=2pt]
\item Un poste par binôme. Pour chaque fait noté : le résultat \emph{et} le site source.
\item Ne recopiez pas des pages entières : deux phrases par domaine suffisent.
\item En fin de rotation, chaque binôme garde sa trouvaille la plus surprenante pour la mise en commun.
\end{itemize}
\subsection*{Domaine 1 --- Histoire : pourquoi le logarithme a-t-il été inventé ?}
\textbf{1.} Qui sont John Napier et Henry Briggs ? Que publie chacun, et quand ?
\par\smallskip
\textbf{2.} Pourquoi les tables de logarithmes ont-elles été imprimées sur des milliers de pages ? Qui les utilisait, et pour quoi faire ?
\par\smallskip
\textbf{3.} Retrouvez la citation attribuée à Laplace sur le logarithme et l'astronome, et expliquez-la en une phrase.
\par\smallskip
{\small\emph{Pistes de recherche :} logarithme Napier 1614 --- Briggs tables 1624 --- règle à calcul --- Laplace astronome}
\par\smallskip
\noindent\makebox[\linewidth]{\dotfill}\par\vspace{1em}
\noindent\makebox[\linewidth]{\dotfill}\par\vspace{1em}
\par\smallskip{\small Sources : \champ{0.96\linewidth}}\par\medskip
\subsection*{Domaine 2 --- Le principe : transformer les produits en sommes}
\textbf{1.} Retrouvez la propriété fondamentale du logarithme : $\log(a \times b) = \,?$ et $\log(a^{n}) = \,?$
\par\smallskip
\textbf{2.} Cherchez comment la règle à calcul utilisait cette propriété, et décrivez-la en deux phrases.
\par\smallskip
{\small\emph{Pistes de recherche :} propriété logarithme du produit --- règle à calcul Gunter Oughtred}
\par\smallskip
\noindent\makebox[\linewidth]{\dotfill}\par\vspace{1em}
\noindent\makebox[\linewidth]{\dotfill}\par\vspace{1em}
\par\smallskip{\small Sources : \champ{0.96\linewidth}}\par\medskip
\subsection*{Domaine 3 --- Lire un exposant sans machine}
\textbf{1.} Retrouvez ce qu'on appelle la caractéristique et la mantisse d'un logarithme décimal.
\par\smallskip
\textbf{2.} Vérifiez en ligne les valeurs de $\log 2$ et $\log 3$ (quatre décimales), puis cherchez comment on calculait ces valeurs avant les calculatrices.
\par\smallskip
{\small\emph{Pistes de recherche :} caractéristique mantisse logarithme --- tables de Briggs 14 décimales}
\par\smallskip
\noindent\makebox[\linewidth]{\dotfill}\par\vspace{1em}
\noindent\makebox[\linewidth]{\dotfill}\par\vspace{1em}
\par\smallskip{\small Sources : \champ{0.96\linewidth}}\par\medskip
\subsection*{Domaine 4 --- Le logarithme dans la vie quotidienne}
\textbf{1.} Rassemblez trois unités du quotidien fondées sur le logarithme (pH, décibels, et une autre de votre choix) et notez leur signification exacte.
\par\smallskip
\textbf{2.} Trouvez l'échelle de magnitude des séismes : que se passe-t-il réellement quand la magnitude augmente de 1 ? De 2 ?
\par\smallskip
{\small\emph{Pistes de recherche :} échelle pH décibel --- magnitude Richter énergie --- séisme magnitude 32}
\par\smallskip
\noindent\makebox[\linewidth]{\dotfill}\par\vspace{1em}
\noindent\makebox[\linewidth]{\dotfill}\par\vspace{1em}
\par\smallskip{\small Sources : \champ{0.96\linewidth}}\par\medskip
\subsection*{Domaine 5 --- Sciences et industrie : compresser les échelles}
\textbf{1.} Retrouvez deux graphiques célèbres tracés en échelle logarithmique (épidémie, astronomie, etc.) et expliquez pourquoi le logarithme y est indispensable.
\par\smallskip
\textbf{2.} Cherchez ce qu'était le papier semi-logarithmique : à quoi servait-il, et qui l'utilise encore ?
\par\smallskip
{\small\emph{Pistes de recherche :} échelle logarithmique épidémie --- papier semi-logarithmique --- ordres de grandeur univers}
\par\smallskip
\noindent\makebox[\linewidth]{\dotfill}\par\vspace{1em}
\noindent\makebox[\linewidth]{\dotfill}\par\vspace{1em}
\par\smallskip{\small Sources : \champ{0.96\linewidth}}\par\medskip
\subsection*{Domaine 6 --- Informatique : chercher en un minimum d'étapes}
\textbf{1.} Cherchez « complexité logarithmique O(log n) » : notez une définition simple et deux algorithmes concernés.
\par\smallskip
\textbf{2.} Trouvez un exemple où une recherche sur un million d'éléments demande moins de vingt étapes, et expliquez pourquoi $2^{20} > 10^{6}$.
\par\smallskip
{\small\emph{Pistes de recherche :} complexité O(log n) --- recherche dichotomique --- algorithme logarithme}
\par\smallskip
\noindent\makebox[\linewidth]{\dotfill}\par\vspace{1em}
\noindent\makebox[\linewidth]{\dotfill}\par\vspace{1em}
\par\smallskip{\small Sources : \champ{0.96\linewidth}}\par\medskip
\newpage
\begin{center}{\Large\bfseries Grille de mise en commun}\\[4pt]{\small Comparer les deux équipes, domaine par domaine}\end{center}
Pour chaque domaine, notez ce que l'équipe papier a établi par le raisonnement, puis ce que l'équipe ordinateur a rapporté du web. Qui a été le plus rapide ? Le plus fiable ? Le plus loin ?\par\medskip
\begin{center}
\begin{tabular}{@{}p{0.3\linewidth}@{\hspace{0.02\linewidth}}p{0.32\linewidth}@{\hspace{0.02\linewidth}}p{0.32\linewidth}@{}}
\toprule
\textbf{Domaine} & \textbf{Équipe papier (raisonnement)} & \textbf{Équipe ordinateur (recherche)} \\
\midrule
{\small\textbf{1.} Histoire : pourquoi le logarithme a-t-il été inventé ?} & \champ{\linewidth} & \champ{\linewidth} \\[18pt]
{\small\textbf{2.} Le principe : transformer les produits en sommes} & \champ{\linewidth} & \champ{\linewidth} \\[18pt]
{\small\textbf{3.} Lire un exposant sans machine} & \champ{\linewidth} & \champ{\linewidth} \\[18pt]
{\small\textbf{4.} Le logarithme dans la vie quotidienne} & \champ{\linewidth} & \champ{\linewidth} \\[18pt]
{\small\textbf{5.} Sciences et industrie : compresser les échelles} & \champ{\linewidth} & \champ{\linewidth} \\[18pt]
{\small\textbf{6.} Informatique : chercher en un minimum d'étapes} & \champ{\linewidth} & \champ{\linewidth} \\[18pt]
\bottomrule
\end{tabular}
\end{center}
\par\medskip{\small Question de synthèse pour la classe : sur quel domaine l'équipe papier a-t-elle été aussi loin que l'équipe ordinateur ? Pourquoi ?}
\end{document}